\input{../tex-inputs/latex-leseansicht-vorspann}

\section[Gustav Schwarzkopf an Arthur Schnitzler, 27. 8. 1929]{L04244 Gustav Schwarzkopf an Arthur Schnitzler, 27. 8. 1929}
\nopagebreak\mylabel{L04244v}
\rehead{ }\normalsize\beginnumbering\briefempfaengerindex{Schnitzler, Arthur@\textsc{Schnitzler, Arthur}!zzzSchwarzkopf, Gustav@\emph{von Gustav Schwarzkopf}!1929-08-271@{27. 8. 1929}|(be}
\toendnotes[C]{\smallbreak\pagebreak[2]}
\correspDesc{Versand  durch Gustav Schwarzkopf am 27. 8. 1929 in Wien
\newline{}Erhalt  durch Arthur Schnitzler im Zeitraum [28. 8. 1929 – 1. 9. 1929?] in Caux}\toendnotes[C]{\smallbreak}
\Standort{CUL, Schnitzler, B 96a.}
\physDesc{Bildpostkarte, 589 Zeichen
\newline{}Handschrift: schwarze Tinte, deutsche Kurrent
\newline{}Versand: Stempel: »\nobreak{}\oindex{I., Innere Stadt@\textbf{I., Innere Stadt}, \emph{Verwaltungsgebiet}|pwk}1 Wien 8, 27. VIII. 29, 20\nobreak{}«.  
\newline{}Schnitzler: mit rotem Buntstift eine Unterstreichung }\toendnotes[C]{\smallbreak}\pstart{}{\pb}I. Tiefer Graben 17. II. 6\oindex{Wien@\textbf{Wien}!I., Innere Stadt@\textbf{I., Innere Stadt}!Tiefer Graben 17@\textbf{Tiefer Graben 17}, \emph{Wohngebäude}|pw}\pend{}\pstart{}Wien\oindex{Wien@\textbf{Wien}, \emph{Verwaltungsgebiet}|pw}.\pend{}{\bigskip}\pstart{}Herrn\pend{}\pstart{}\textsc{D\textsuperscript{r}  Arthur
                  Schnitzler}\pend{}\pstart{}\textsc{Caux (Schweiz)\oindex{Caux@\textbf{Caux}|pw}}\pend{}\pstart{}\textsc{\label{K_L04244-1v}\edtext{Palace-Hôtel\oindex{Caux Palace Hotel@\textbf{Caux Palace Hotel}, \emph{Hotel}|pw}}{\lemma{\textnormal{\emph{Palace-Hôtel}}}\Cendnote{\textnormal{Schnitzler
                        und Pollaczek\pwindex{Pollaczek, Clara Katharina 15.\,1.\,1875 Wien – 22.\,7.\,1951 ebd.@\textsc{Pollaczek, Clara Katharina} (15.\,1.\,1875 Wien – 22.\,7.\,1951 ebd.), \emph{Schriftstellerin}|pwk} wohnen im Grand Hotel Excelsior Regina\oindex{Grand Hotel Excelsior Regina@\textbf{Grand Hotel Excelsior Regina}, \emph{Hotel}|pwk}.}}}\label{K_L04244-1}}\pend{}{\bigskip}
\pstart
           \noindent{}\centering{}{\pb}\textcolor{gray}{\textbf{Badgastein, Salzburg}}\oindex{Bad Gastein@\textbf{Bad Gastein}, \emph{Hauptstadt}|pw}\pend
           \vspace{1em}
\pstart
           \raggedleft{}{\pb}27. VIII. 29\pend
           
\pstart{}Lieber Arthur!\pend\vspace{0.5em}
\pstart
           Dank für ihre \label{K_L04244-2v}\edtext{Karte}{\lemma{\textnormal{\emph{Karte}}}\Cendnote{\textnormal{{XXXX ref} XXXX}}}\label{K_L04244-2}. Freue mich, daß Sie es{ }ſo gut getroffen haben, der Aufenthalt,{ }ſoll ja, wie man mir hier{ }ſagt, zu den{ }ſchönſten zählen. Ich bin Montag
               Abend (19.) zurück gekommen, war gut aufgehoben und habe nichts
               Unangenehmes erlebt. Hier wieder die übliche Einteilung: Volksgarten\oindex{Wien@\textbf{Wien}!I., Innere Stadt@\textbf{I., Innere Stadt}!Volksgarten@\textbf{Volksgarten}, \emph{Park}|pw}, Rathauspark\oindex{Wien@\textbf{Wien}!I., Innere Stadt@\textbf{I., Innere Stadt}!Rathauspark@\textbf{Rathauspark}, \emph{Park}|pw}, Burggarten\oindex{Wien@\textbf{Wien}!I., Innere Stadt@\textbf{I., Innere Stadt}!Burggarten@\textbf{Burggarten}, \emph{Park}|pw}. \textcolor{gray}{Neues} ka{\geminationn}
               ich Ihnen nicht melden, die meiſten Beka{\geminationn}ten{ }ſind noch
               abweſend. Ich grüße Sie {\pb}herzlichſt und
               bitte Sie auch an Frau P.\pwindex{Pollaczek, Clara Katharina 15.\,1.\,1875 Wien – 22.\,7.\,1951 ebd.@\textsc{Pollaczek, Clara Katharina} (15.\,1.\,1875 Wien – 22.\,7.\,1951 ebd.), \emph{Schriftstellerin}|pw} meine allerſchönsten Grüße zu
               beſtellen\pend
           
\pstart
           Ihr{\\[\baselineskip]}\spacefill\mbox{Gustav}\pend
           \leftskip=0em{}\selectlanguage{ngerman}\endnumbering\briefempfaengerindex{Schnitzler, Arthur@\textsc{Schnitzler, Arthur}!zzzSchwarzkopf, Gustav@\emph{von Gustav Schwarzkopf}!1929-08-271@{27. 8. 1929}|)be}\mylabel{L04244h}
\begin{anhang}
\end{anhang}\newcommand{\dateiname}{L04244}\newcommand{\titel}{Gustav Schwarzkopf an Arthur Schnitzler, 27. 8. 1929}\newcommand{\editorInnen}{Herausgegeben von Jahnke, SelmaMüller, Martin Anton}\input{../tex-inputs/latex-leseansicht-abspann}
